\section{Experiments}
\subsection{Evaluation Tasks.}
We evaluate our approach on long-horizon manipulation tasks that require retrieving diverse information types---spatial, relational, numerical, and event-time---from past interactions, including human-robot collaborative scenarios.
These tasks span five scenarios reflecting diverse memory demands often faced in household settings (Fig.~\ref{fig:rw_task_viz}):
\vspace{0.5em}
\\
\textbf{Retrieve Object} (\textit{Spatial information}).
This task requires recalling object locations observed earlier in the episode; the policy must fetch a granola bar from a location previously observed on the kitchen countertop, which is no longer visible at the time of action.
Across episodes, the object’s initial location varies between the left and right sides of the countertop.
% \vspace{0.5em}
\\
\textbf{Return to Same Container} (\textit{Object relations}).
This task requires recalling earlier relations among objects; after placing a bowl in a sink, the policy must later retrieve the plate on which the bowl was originally placed.
Across episodes, the involved objects and their relations vary.
% \vspace{0.5em}
\\
\textbf{Store $N$ objects} (\textit{Numerical information}).
This task requires recalling a numerical count over time; the robot should place $n$ chip packets in a drawer before closing it, with the packets becoming occluded once placed.
Across episodes, both the initial number of chip packets and the target count $n$ vary.
% \vspace{0.5em}
\\
\textbf{Heat Stove} (\textit{Event-time information}).
This task requires recalling when an event occurred to act after a specified duration; the robot must heat a stove for a fixed duration.
Across episodes, the required heating duration varies between $4$, $6$, and $8$ minutes.
% \vspace{0.5em}
\\
\textbf{Store $N$ objects} (\textit{Spatial information, Human-robot collaboration}).
This task requires tracking a numerical count in a collaborative setting: a human places some purple cups first, then the robot completes the total of three purple cups in the shopping bag, and then adds a green cup to it.
Across episodes, the number of initial purple cups in the scene ranges from $3$ to $4$, and the human places $1$ to $3$ cups.
% \textbf{Method} & \multicolumn{4}{c}{\textbf{Information Type}} \\
% \cmidrule(lr){2-5}
%  & \textbf{Retrieve Object} & \textbf{Return to Same Container} & \textbf{Store N Objects} & \textbf{Heat Stove for $T$ mins.} \\
% \begin{table*}[t]
% \centering
% \small
% \setlength{\tabcolsep}{7pt}
% \renewcommand{\arraystretch}{1.2}
% \begin{threeparttable}
% \begin{tabular}{lcccc}
% \toprule
% \textbf{Method} 
% & \textbf{Retrieve Object} 
% & \textbf{Return to Same Container} 
% & \textbf{Store N Objects} 
% & \textbf{Heat Stove for $T$ mins.} \\
% 
% & \textit{(Spatial Info.)} 
% & \textit{(Relational Info.)} 
% & \textit{(Numerical Info.)} 
% & \textit{(Event-time Info.)} \\
% \midrule
% LSTM              & $0.14$ & $0.12$ & -- & -- \\
% Mamba             & $0.20$ & $0.18$ & -- & -- \\
% TransformerXL     & $0.12$ & $0.20$ & -- & -- \\
% \midrule
% Local Attention        & $0.13$    & $0.16$    & $0.18$ & $0.11$ \\
% Strided Attention      & $0.20$ & $0.28$ & -- & -- \\
% Hierarchical Attention & --     & --     & -- & -- \\
% Gated Attention & --     & --     & -- & -- \\
% \midrule
% Standard Transformer     & $0.26 \pm 0.02$ & $0.23 \pm 0.04$ & $0.12 \pm 0.01$ & $0.27\pm0.09$ \\
% Scene Memory Transformer & --     & --     & -- & -- \\
% SAM2Act++                & $0.39$ & $0.11$ & $\mathbf{0.27}$ & $0.04$ \\
% ReMemBer                 & $0.36$ & $0.13$ & $0.21$  & $0.00$ \\
% Hand-Designed Features   & $\mathbf{0.68}$    & $0.15$ & $0.21$ & $0.13$ \\
% \midrule
% HALO w/o VQA         & $0.42 \pm 0.07$ & $0.29 \pm 0.06$ & $0.17\pm0.04$ & $0.37 \pm 0.11$ \\
% HALO                 & $0.64 \pm 0.02$ & $\mathbf{0.32 \pm 0.02}$ & $0.26 \pm 0.03$ & $\mathbf{0.40 \pm 0.11}$ \\
% \bottomrule
% \end{tabular}
% \caption{\rutav{Change from information type to task names}}
% \end{threeparttable}
% \end{table*}

\begin{table*}[t]
\centering
\caption{Evaluation on simulation benchmark with $50$ rollouts per task.}
\vspace{-8pt}
\small
\setlength{\tabcolsep}{7pt}
\renewcommand{\arraystretch}{1.2}
\begin{threeparttable}

\resizebox{\textwidth}{!}{%
\label{tab:sim_all}
\begin{tabular}{lccccc}
\toprule
\textbf{Method} 
& \textbf{Retrieve Object} 
& \textbf{Return to Same Container} 
& \textbf{Store N Objects} 
& \textbf{Heat Stove for $T$ mins.}
& \textbf{Average} \\

& \textit{(Spatial Info.)} 
& \textit{(Relational Info.)} 
& \textit{(Numerical Info.)} 
& \textit{(Event-time Info.)}
& \\
\midrule
Standard Transformer     
& $0.26$ & $0.23$ & $0.12$ & $0.27$ & $0.22$ \\
Scene Memory Transformer & $0.53$ & $0.25$ & $0.17$ & $\mathbf{0.40}$ & $0.34$ \\
SAM2Act++                & $0.39$ & $0.11$ & $\mathbf{0.27}$ & $0.04$ & $0.20$ \\
ReMemBer                 & $0.36$ & $0.13$ & $0.21$  & $0.00$ & $0.18$ \\
Hand-Designed Features   & $\mathbf{0.68}$ & $0.15$ & $0.21$ & $0.13$ & $0.29$ \\
Token Merging            &  $0.29$ &  $0.24$ & $0.25$ &  $0.38$ & $0.29$ \\
\midrule
% HALO w/o VQA         & $0.42 \pm 0.07$ & $0.29 \pm 0.06$ & $0.17\pm0.04$ & $0.37 \pm 0.11$ & $0.31$ \\
% HALO                 & $0.64 \pm 0.02$ & $\mathbf{0.32 \pm 0.02}$ & $0.26 \pm 0.03$ & $\mathbf{0.40 \pm 0.11}$ & $\mathbf{0.41}$ \\
HALO w/o VQA         & $0.42$ & $0.29$ & $0.17$ & $0.37$ & $0.31$ \\
HALO                 & $0.64$ & $\mathbf{0.32}$ & $0.26$ & $\mathbf{0.40}$ & $\mathbf{0.41}$ \\
\bottomrule
\end{tabular}%
}
\end{threeparttable}
\end{table*}

% \begin{table*}[t]
% \centering
% \small
% \setlength{\tabcolsep}{4pt}
% \renewcommand{\arraystretch}{1.2}
% \caption{Evaluation on real-world tasks with $20$ rollouts per task}
% \begin{threeparttable}
% \label{tab:rw_all}
% \begin{tabular}{lcccccc}
% \toprule
% \textbf{Method} 
% & \textbf{Retrieve Obj.} 
% & \textbf{Return to Same Container} 
% & \textbf{Store N Objs.} 
% & \textbf{Heat Stove for $T$ mins.}
% & \textbf{Store N Objs.}
% & \textbf{Average} \\
% 
% & \textit{(Spatial Info.)} 
% & \textit{(Relational Info.)} 
% & \textit{(Numerical Info.)} 
% & \textit{(Event-time Info.)}
% & \textit{(Spatial info., Human-Robot Collab.)}
% & \\
% \midrule
% Standard Transformer & $0.40$ & $0.30$ & $0.20$ & $0.40$ & $0.50$ & $0.36$ \\
% HALO                 & $\mathbf{0.55}$ & $\mathbf{0.40}$ & $\mathbf{0.55}$ & $\mathbf{0.60}$ & $\mathbf{0.65}$ & $\mathbf{0.55}$ \\
% \bottomrule
% \end{tabular}
% \end{threeparttable}
% \end{table*}

\begin{table*}[t]
\centering
\small
\setlength{\tabcolsep}{5pt}
\renewcommand{\arraystretch}{1.2}
\vspace{-12pt}
\caption{Evaluation on real-world tasks with $20$ rollouts per task}
\vspace{-16pt}
\begin{threeparttable}
\label{tab:rw_all}
\begin{tabular}{lcccccc}
\toprule
\textbf{Method} 
& \textbf{Retrieve Obj.}
& \textbf{Return to Container}
& \textbf{Store N Objs.}
& \textbf{Heat Stove}
& \textbf{Store N Objs.}
& \textbf{Average} \\
& 
& 
& 
& \textbf{for $T$ mins.}
& Human-Robot
& \\
& \textit{(Spatial Info.)}
& \textit{(Relational Info.)}
& \textit{(Numerical Info.)}
& \textit{(Event-time Info.)}
& \textit{Spatial Info.}
& \\
\midrule
Standard Transformer & $0.40$ & $0.30$ & $0.20$ & $0.40$ & $0.50$ & $0.36$ \\
HALO                 & $\mathbf{0.55}$ & $\mathbf{0.40}$ & $\mathbf{0.55}$ & $\mathbf{0.60}$ & $\mathbf{0.65}$ & $\mathbf{0.55}$ \\
\bottomrule
\end{tabular}
\end{threeparttable}
\vspace{-10pt}
\end{table*}
\subsection{Baselines.}
With the goal of reducing memory size or introducing inductive priors about what should be stored, prior work proposes several alternatives:
\vspace{0.5em}
\\
\textbf{SAM2Act++}~\cite{sam2act} uses object-centric memory by storing object center pixel coordinates.
This induces a strong spatial prior and reduces memory size by storing only object locations.
% \vspace{0.5em}
\\
\textbf{ReMemBer}~\cite{anwar2025remembr} stores text summaries of past observations generated by a VLM.
This introduces semantic priors while compressing history into text.
% \vspace{0.5em}
\\
\textbf{Hand-designed Features} uses human-designed rules to select which task-relevant features to store or discard, similar to prior work MemER~\cite{memer}.
It encodes human priors about what information is important while keeping memory compact.
% \vspace{0.5em}
\\
\textbf{Scene Memory Transformer (SMT)}~\cite{fang2019scene} compresses the entire history into a single scene-level embedding via attention.
This reduces memory footprint but forces all past information into a fixed-size representation.
% \vspace{0.5em}
\\
\textbf{Token Merging}~\cite{memoryvla} compresses the history by merging tokens with similar embeddings that are temporally adjacent to maintain a fixed budget of memory.

\subsection{Results}
\textit{How well does \ourmethod{} perform compared to prior visuomotor imitation learning methods that incorporate alternative inductive priors?} (Table ~\ref{tab:sim_all})
\vspace{0.5em}
\\
\textbf{Comparison to object-centric memory.}
Compared to SAM2Act++, \ourmethod{} improves average task success by $21\%$ points.
SAM2Act++ performs well when tasks align with its object-centric prior, such as counting objects, but struggles with tasks that require non-spatial memory, particularly recalling event-time information.
This highlights the limitation of relying on a single type of prior.
\\
\textbf{Comparison to text-based summarization.}
\ourmethod{} outperforms ReMemBer by $23\%$ points on average.
Text summaries generated by VLM often capture coarse semantic information, but may omit details needed for control.
Textual summaries may capture object locations but omit action information for navigating to the object, explaining poor `Retrieve Object' performance.
It suggests that while VLM priors are useful, text-based compression alone is insufficient for control.
\\
\textbf{Comparison to hand-designed rule-based filtering.}
Compared to hand-designed features, \ourmethod{} achieves an absolute improvement of $12\%$.
While this performs on par with \ourmethod{} on spatial tasks, it remains less competitive than \ourmethod{} on other tasks.
This suggests that learning what to retrieve directly from data using~\ourmethod{} not only removes manually designed task-specific priors but also improves performance, possibly by enabling the policy to exploit additional information for decision-making (App. Sec.~\ref{sec:app:retrieved_info_vis}).
\\
\textbf{Comparison to compressed memory representations.}
We compare \ourmethod{} against two approaches that compress history: Scene Memory Transformer (SMT), which uses a learned pooled attention layer to compress history into a single embedding, and Token Merging, which merges nearby tokens to maintain a fixed memory budget.
\ourmethod{} outperforms SMT by $7\%$ and Token Merging by $12\%$ points.
While compression reduces redundancies and noise in memory, it can lose granular information necessary for low-level action prediction~\cite{shao2025tokens}.
In contrast, \ourmethod{} stores all information in memory, enabling access to fine-grained details and performing selective retrieval to reduce the impact of noise.
% In conclusion, in contrast to prior work on different types of compression or inductive biases for selecting what should be stored,~\ourmethod{} stores all information in the history, maintains granular information, and performs selective retrieval.

Overall, we observe that no single method wins across all tasks. However,~\ourmethod{} remains competitive without task-specific assumptions or hand-designed rules, making it versatile with less engineering effort across information types. This versatility is reflected in average performance across tasks, alongside competitive performance on individual ones.
\\
\\
\textit{How effective is \ourmethod{} in real-world tasks?} (Table ~\ref{tab:rw_all})
% \vspace{0.5em}
% \\

We observe a similar trend in real-world settings, where \ourmethod{} consistently outperforms the standard Transformer baseline by $19\%$. Notably, this improvement persists on challenging tasks, including storing multiple objects, tasks completed by the human partner, and long-horizon tasks like heating a pot for up to 8 minutes, which demand diverse types of information and long-context memory.

In addition, we measure manipulation and memory failures in real-world evaluations, finding that HALO reduces them by $8\%$ and $25\%$ absolute over full attention in the `Retrieve Object' task, respectively. These results support our hypothesis that HALO reduces model drift (fewer manipulation failures) and improves memory retrieval (fewer memory failures).
\\
\\
\begin{table}[t]
\centering
\caption{Evaluation on simulation with $50$ rollouts per task.}
\vspace{-8pt}
\small
\setlength{\tabcolsep}{6pt}
\renewcommand{\arraystretch}{1.2}
\begin{threeparttable}
\resizebox{0.95\columnwidth}{!}{%
\begin{tabular}{lcc}
\toprule
\textbf{Method}
& \textbf{Retrieve Object} 
& \textbf{Return to Container} \\
% & \textit{(Spatial)} 
% & \textit{(Relational)} \\
\midrule
LSTM          & $0.14$ & $0.12$ \\
Mamba         & $0.20$ & $0.18$ \\
TransformerXL & $0.12$ & $0.20$ \\
\midrule
Window Attention   & $0.13$ & $0.16$ \\
Strided Attention & $0.20$ & $0.28$ \\
Hierarchical Attention & $0.28$ & $0.22$ \\
Gated Attention        & $0.45$ & $0.28$ \\
\midrule
HALO & $\mathbf{0.64}$ & $\mathbf{0.32}$ \\
\bottomrule
\end{tabular}
}
\end{threeparttable}
\vspace{-15pt}
\end{table}
% LSTM~\cite{lstm}          & $0.14$ & $0.12$ \\
% Mamba~\cite{mamba}         & $0.20$ & $0.18$ \\
% TransformerXL~\cite{dai2019transformer} & $0.12$ & $0.20$ \\
% \midrule
% Local Attention~\cite{vaswani2017attention}   & $0.13$ & $0.16$ \\
% Strided Attention~\cite{child2019generating} & $0.20$ & $0.28$ \\
% Hierarchical Attention~\cite{hierarchicalattention} & $0.28$ & $0.22$ \\
% Gated Attention~\cite{gatedattention}        & $0.45$ & $0.28$ \\
\textit{What are the effects of VLM-induced priors and Top-$K$ sparsification on performance?} (Table ~\ref{tab:sim_all})
\vspace{0.5em}
\\
\textbf{Effect of VLM-induced priors.}
We compare \ourmethod{} against a variant trained without VQA supervision. Adding VLM-induced priors improves average success by $10\%$, suggesting that VQA supervision helps the policy identify relevant history, particularly in tasks that require recalling past object placements or numerical information. Without it, the retrieval mechanism may often attend to irrelevant history, leading to poor performance at test time.
% \vspace{0.5em}
\\
\textbf{Effect of Top-$K$ sparsification.}
We compare Top-$K$ retrieval to full attention over memory.
Top-$K$ improves performance by $9\%$.
Restricting retrieval reduces the amount of irrelevant or noisy information incorporated into decision-making, which is particularly important in long-horizon tasks where only a few past events are relevant to the current action.

See App. Sec.~\ref{sec:app:spurious_vis} and Sec.~\ref{sec:app:model_drift} for more analysis.
\\
\\
\textit{How does \ourmethod{} compare to alternative sparsification and memory mechanisms?}
\vspace{0.5em}
\\
\textbf{Comparison to alternative sparsification methods.}
We compare Top-$K$ sparsification to window attention~\cite{child2019generating}, strided attention~\cite{child2019generating}, hierarchical attention~\cite{beltagy2020longformer}, and gated attention~\cite{gatedattention}.
\ourmethod{} outperforms the best alternative, gated attention, by $11\%$.
The sparsification methods relying on fixed access patterns may miss task-relevant events occurring at unpredictable times.
In contrast, Top-$K$ selects memory entries based on content relevance, allowing the policy to retrieve information tied to the task rather than position.
% \vspace{0.5em}
\\
\textbf{Comparison to non-associative memory mechanisms.}
We compare against models that compress history into latent states, including LSTM~\cite{lstm} ($-20\%$), Mamba~\cite{mamba} ($-14\%$), and Transformer-XL~\cite{dai2019transformer} ($-12\%$).
While compressed-state models summarize past information efficiently, they may discard details needed for recalling specific past events.
Selective retrieval from explicit memory allows the policy to access such details when required.
\\
\\
\textit{What are the effects of different design choices in~\ourmethod{}?}
\vspace{0.5em}
\\
\textbf{Number of retrieved entries ($k$).}
% \begin{table}[h]
\centering
\vspace{-10pt}
% \caption{Effect of the number of retrieved entries ($k$).}
% \vspace{-10pt}
\small
\setlength{\tabcolsep}{6pt}
\renewcommand{\arraystretch}{1.1}
\resizebox{0.9\columnwidth}{!}{%
\begin{tabular}{lccccc}
\toprule
\textbf{$k$} & 4 & 8 & 12 & 16 & 24 \\
\midrule
\textbf{Retrieve Object} 
& $0.40$
& $\mathbf{0.52}$
& $0.48$
& $0.33$
& $0.44$ \\
\bottomrule
\end{tabular}%
}
\end{table}
We ablate the number of retrieved history entries ($k$) used by the attention mechanism. A moderate value ($k=8$) achieves the best performance ($52\%$ success). Smaller $k$ (\textit{e.g.}, $k=4$, $40\%$ success) omits useful context, while larger $k$ incorporates irrelevant history ($48\%$ for $k=12$, $33\%$ for $k=16$, $44\%$ for $k=24$). We fix $k=8$ across all tasks. Developing adaptive strategies that retrieve only the necessary amount of information at each step is a promising direction for future work.
% \vspace{0.5em}
\\
\textbf{Co-training vs.\ pretraining.}
% \begin{table}[t]
\centering
\caption{effect of co-training on task success rate.}
\small
\resizebox{0.9\columnwidth}{!}{%
\begin{tabular}{lc}
\toprule
\textbf{training strategy} & \textbf{retrieve object} \\
& \textit{(spatial info.)} \\
\midrule
no vqa (action only)             & $0.42$ \\
vqa pretrain + action finetuning & $0.44$ \\
vqa + action co-training         & $\mathbf{0.64}$ \\
\bottomrule
\end{tabular}
}
\end{table}
We compare joint co-training of VQA and action prediction against a two-stage approach (VQA pretraining followed by action finetuning). Co-training VQA and action prediction achieves $64\%$ success, outperforming pretrain-then-finetune ($44\%$) and no-VQA training ($42\%$) by $20$ and $22$ points, respectively. The minimal gain from separate pretrain-then-finetune compared to no-VQA suggests that VQA knowledge is lost during fine-tuning, whereas co-training effectively shapes retrieval toward task-relevant information.
% \vspace{0.5em}
\\
\textbf{Stage-wise VQA generation vs. single video prompt generation.}
We compare HALO's stage-wise pipeline against a single-prompt baseline that generates QA pairs directly from long videos.
HALO achieves $\mathbf{60\%}$ points more accurate QAs and $\mathbf{20\%}$ points fewer hallucinations.
The results suggest SOTA VLMs struggle with direct long-video reasoning for question-answer generation, highlighting the importance of the proposed stage-wise pipeline for generating high-quality VQA (Examples in App. Sec.~\ref {sec:app:vqa_viz}).
% \vspace{0.5em}
% \\
% \textbf{Stage-wise VQA generation vs. single video prompt generation}
% HALO uses a stage-wise pipeline: (i) LLM and SAM extract per-frame granular task-relevant information, (ii) activity summarization captures subtask-level information, and (iii) both are combined to produce VQA pairs.
% Empirically, this yields $\mathbf{+60\%}$ pts. more accurate QAs and $\mathbf{-20\%}$ pts. fewer hallucinations compared to a single-prompt baseline.
% This result suggests that SOTA VLMs often produce question-answer pairs directly from long videos, often hallucinating objects or generating incorrect answers.